\chapter*{از داده تا جهان فیزیکی: نتیجه‌ی بحث}
\addcontentsline{toc}{chapter}{از داده تا جهان فیزیکی: نتیجه‌ی بحث}
\markboth{از داده تا جهان فیزیکی}{از داده تا جهان فیزیکی}

در آغاز این کتابچه، «داده» واژه‌ای آشنا به نظر می‌رسید؛ چیزی شبیه عددی در یک جدول، رشته‌ای در حافظه یا پاسخی که یک سامانه ثبت می‌کند. اما هرچه جلوتر آمدیم، روشن شد که این تصویر ساده چند پرسش را پنهان می‌کند: چه چیزی باعث می‌شود یک تفاوت را ثبت کنیم؟ چه زمانی آن تفاوت برای ما اطلاعات می‌شود؟ و چگونه چیزی که در ذهن یا در یک دستگاه به شکل بازنمایی حضور دارد، در جهان واقعی قابل نگهداری و انتقال می‌شود؟

فصل نخست میان داده، اطلاعات، دانش و خرد فاصله گذاشت. داده به‌تنهایی تضمین نمی‌کند که چیزی فهمیده شده باشد؛ برای رسیدن از یک ثبت خام به اطلاعات، زمینه و تفسیر لازم است. فصل دوم نشان داد که بازنمایی نیز صرفاً شباهت ظاهری یا قراردادی دلخواه نیست. چیزی می‌تواند نماینده‌ی چیز دیگری باشد، اما این رابطه به زمینه، سامانه‌ی تفسیر و نوع پیوند میان آن دو وابسته است.

در فصل سوم، پرسش از بازنمایی به سیگنال فیزیکی رسید. یک پیام با کاغذ، موج صوتی یا وضعیت‌های یک مدار یکی نیست؛ بااین‌حال، بدون یک حامل قابل تشخیص نمی‌تواند از راهی مادی ثبت یا منتقل شود. سیگنال همان معنا نیست، بلکه تغییری فیزیکی است که در یک سامانه‌ی مشخص می‌تواند حامل پیام باشد. از همین‌جا اهمیت نویز، احتمال، ظرفیت کانال، نمونه‌برداری و آستانه‌های منطقی آشکار می‌شود: سامانه باید بتواند تفاوت مورد نظر را از تغییرهای ناخواسته جدا کند.

این تمایزها کمک می‌کنند فرمول‌ها را هم در جای درستشان ببینیم. آنتروپی شانون عدم‌قطعیت یک منبع احتمالاتی را اندازه می‌گیرد، نه معنای انسانی پیام را. پیچیدگی کولموگروف به کوتاه‌ترین توصیف یک رشته مربوط است، نه به میزان فهم یا اهمیت آن. اصل لاندائر نیز حد پایین هزینه‌ی گرماییِ عملیات منطقیِ برگشت‌ناپذیرِ پاک‌کردن اطلاعات را بیان می‌کند؛ نه اینکه هر محاسبه یا هر جابه‌جایی اطلاعات الزاماً همان هزینه را داشته باشد.

وقتی بیت در یک مدار تحقق پیدا می‌کند، دیگر با یک صفر و یکِ جدا از جهان روبه‌رو نیستیم. با ولتاژ، بار، ماده، نویز، گرما و محدودیت‌های یک سامانه‌ی واقعی روبه‌رو هستیم. این واقعیت از ارزش بازنمایی انتزاعی کم نمی‌کند؛ فقط یادآوری می‌کند که هر بازنماییِ قابل استفاده باید جایی تحقق پیدا کند و راهی برای تشخیص و تفسیر داشته باشد.

پس پاسخ این کتابچه یک شعار ساده مانند «همه‌چیز فیزیکی است» نیست. نکته‌ی دقیق‌تر این است: برای فهم داده باید هم رابطه‌ی معنایی و تفسیری آن را بشناسیم و هم شرایط مادی‌ای را که ثبت، تمایز، انتقال و نگهداری آن را ممکن می‌کنند. هرگاه با داده‌ای روبه‌رو شدیم، پرسش فقط این نیست که «چه چیزی نوشته شده؟»؛ باید بپرسیم کدام تفاوت ثبت شده، در چه بستری پایدار مانده، با چه قراردادی رمزگشایی می‌شود و برای چه سامانه‌ای چه چیزی را بازنمایی می‌کند.

و شاید پرسش نهایی همین باشد: در یک سامانه، دقیقاً از چه لحظه‌ای یک وضعیت فیزیکی صرفاً یک وضعیت نیست و به داده‌ای تبدیل می‌شود که می‌توان آن را خواند، منتقل کرد و درباره‌اش استدلال کرد؟
